\documentclass[11pt]{article}
\usepackage[utf8]{inputenc}
\usepackage{amsmath,amssymb,amsthm}
\usepackage[margin=1in]{geometry}

\newtheorem{lemma}{Lemma}
\newtheorem{theorem}{Theorem}

\title{Problem 788: Dominating Numbers}
\author{}
\date{}

\begin{document}
\maketitle

\section{Problem Statement}

A positive integer with exactly $n$ digits is called dominating if some digit appears more than $n/2$ times. Let $D(N)$ be the number of dominating positive integers with at most $N$ digits. Compute
\[
  D(2022) \bmod 1{,}000{,}000{,}007.
\]

\section{Mathematical Foundation}

For a fixed length $n$, define
\[
  T(n) = \sum_{k=\lfloor n/2 \rfloor + 1}^{n} \binom{n}{k} 9^{n-k}.
\]

\begin{lemma}
For fixed $n$ and fixed digit $d$, the number of length-$n$ strings in which $d$ appears exactly $k$ times is
\[
  \binom{n}{k}9^{n-k}.
\]
\end{lemma}

\begin{proof}
Choose the $k$ positions occupied by $d$, then fill the remaining positions with any of the $9$ digits different from $d$.
\end{proof}

\begin{lemma}
For a fixed length $n$, at most one digit can appear more than $n/2$ times.
\end{lemma}

\begin{proof}
If two different digits each appeared more than $n/2$ times, the total number of digit occurrences would exceed $n$.
\end{proof}

\begin{theorem}
The number of dominating positive integers with exactly $n$ digits is
\[
  9T(n).
\]
\end{theorem}

\begin{proof}
By the previous lemma there is no overlap between different dominating digits, so we may count by the dominating digit.

For a fixed nonzero digit $d \in \{1,\dots,9\}$ and fixed $k>n/2$:
\begin{itemize}
  \item total strings with exactly $k$ copies of $d$: $\binom{n}{k}9^{n-k}$,
  \item among them, those with leading zero: $\binom{n-1}{k}9^{n-1-k}$.
\end{itemize}
So the count for $d \neq 0$ is
\[
  \binom{n}{k}9^{n-k} - \binom{n-1}{k}9^{n-1-k}.
\]
For $d=0$, the leading digit cannot be zero, so the count is
\[
  \binom{n-1}{k}9^{n-k}.
\]
Summing over $d=1,\dots,9$ and then adding the $d=0$ contribution gives
\[
  9\left(\binom{n}{k}9^{n-k} - \binom{n-1}{k}9^{n-1-k}\right)
  + \binom{n-1}{k}9^{n-k}
  = 9\binom{n}{k}9^{n-k}.
\]
Now sum over all $k>n/2$.
\end{proof}

Therefore
\[
  D(N) = 9 \sum_{n=1}^{N} \sum_{k=\lfloor n/2 \rfloor + 1}^{n} \binom{n}{k} 9^{n-k}.
\]

\section{Algorithm}

\begin{enumerate}
  \item Precompute factorials and inverse factorials modulo $10^9+7$ up to $N=2022$.
  \item Precompute powers of $9$.
  \item For each $n=1,\dots,2022$, evaluate
  \[
    T(n) = \sum_{k=\lfloor n/2 \rfloor + 1}^{n} \binom{n}{k} 9^{n-k},
  \]
  then add $9T(n)$ to the total.
\end{enumerate}

\section{Complexity Analysis}

\begin{itemize}
  \item Time: $O(N^2)$.
  \item Space: $O(N)$.
\end{itemize}

\section{Answer}

\[
\boxed{471745499}
\]

\end{document}
